\subsection{Open-World Sensing and Perception}
\label{sec:rw_sensing}
% =============================================================

In closed-world settings, perception systems are designed to recognize a fixed set of categories, objects, or states defined at training time. Open-world perception, by contrast, requires the ability to detect and reason about entities and situations that were not anticipated during system design. Early work on open set recognition \citep{scheirer2013toward, bendale2016towards} extended standard classifiers to reject inputs that do not belong to any known category, typically through confidence thresholding or distribution boundary modeling. Out-of-distribution detection methods \citep{hendrycks2017baseline} further developed principled approaches for identifying inputs that fall outside the training distribution.

More recently, vision-language models (VLMs) have transformed the landscape of open-world perception. CLIP \citep{radford2021learning} demonstrated that joint vision-language pretraining enables zero-shot recognition of arbitrary categories described in natural language, fundamentally relaxing the closed-vocabulary constraint. This capability has been extended to object detection through open-vocabulary detectors such as ViLD \citep{gu2022openvocabulary} and OWL-ViT \citep{minderer2022simple}, which can localize objects described by free-form text queries without category-specific training. Foundation models for segmentation, such as the Segment Anything Model \citep{kirillov2023segment}, further demonstrate that general-purpose visual representations can support open-ended perception tasks.

Despite these advances, a critical gap remains between perceptual novelty detection and the needs of an acting agent. Recognizing that an object is unfamiliar or that a scene has changed is necessary but not sufficient: the agent must also determine whether the novelty is functionally relevant to its current task and how it should modify its behavior in response. Current perception systems largely operate as passive classifiers and do not reason about the implications of detected novelty for downstream planning and execution.